%!TEX root = ../main.tex

\chapter{Requisitos del sistema}
\label{ch:requisitos}

A partir del estado del arte presentado en el \cref{ch:estado_del_arte} y del
alcance definido en la \cref{sec:alcance}, en este capítulo se definen los
requisitos que el sistema debe cumplir. Los requisitos se organizan en cuatro
categorías: funcionales (RF), de rendimiento (RP), restricciones del proyecto
(RC) y normativos (RN). Para justificar estos requisitos, en cada requisito se incluye una o más referencias, ya sean normas aeronáuticas, alguna sección del estado del arte, o una decisión del proyecto; las cuales apoyan la decisión tomada para ese requisito.

Es importante recordar que para los requisitos de rendimiento, como no se pretende llegar a certificación aeronáutica, los umbrales definidos son solo orientativos, son coherentes con la norma, pero no son de cumplimiento estricto. Dicho esto, los valores son definidos para disponer de criterios objetivos a los que hacer referencia en las secciones de validación.


\section{Requisitos funcionales}
\label{sec:req_funcionales}

Los requisitos funcionales describen las capacidades que el sistema debe
proporcionar. Se derivan principalmente de la disposición estándar Basic-T del PFD definida en \cite{ac25_11b}, complementada con funcionalidades adicionales añadidas para reforzar la finalidad educativa de la plataforma.

\begin{table}[H]
    \centering
    \caption{Requisitos funcionales del sistema}
    \label{tab:req_funcionales}
    \begin{tabularx}{\linewidth}{l X l}
        \toprule
        \textbf{ID} & \textbf{Descripción} & \textbf{Fuente} \\
        \midrule
        RF01 & Estimar la actitud (cabeceo y alabeo) de la plataforma en tiempo real. & \cite{ac25_11b} \\
        RF02 & Estimar el rumbo magnético de la plataforma en tiempo real. & \cite{ac25_11b} \\
        RF03 & Calcular la altitud respecto al nivel del mar a partir de la presión atmosférica. & \cite{icao_isa}, \cite{ac25_11b} \\
        RF04 & Estimar la velocidad vertical (tasa de ascenso y descenso). & \cite{ac25_11b} \\
        RF05 & Estimar la velocidad respecto al suelo a partir de datos GNSS. & Decisión de proyecto \\
        RF06 & Medir la distancia al suelo mediante un sensor de proximidad principal, complementado con un segundo sensor de tecnología distinta con fines comparativos. & Decisión de proyecto \\
        RF07 & Transmitir la telemetría de forma inalámbrica desde el nodo sensor hasta un receptor. & \cref{sec:ea_contribucion} \\
        RF08 & Codificar la telemetría conforme al formato lógico del protocolo ARINC 429. & \cite{arinc429}, \cref{sec:ea_arinc} \\
        RF09 & Visualizar los datos calculados en un PFD con disposición Basic-T. & \cite{ac25_11b} \\
        \bottomrule
    \end{tabularx}
\end{table}


\section{Requisitos de rendimiento}
\label{sec:req_rendimiento}

Los requisitos de rendimiento cuantifican el comportamiento esperado del sistema.
Los valores se derivan de los estándares de pantallas electrónicas de vuelo, complementados con valores obtenidos de plataformas AHRS con sensores MEMS de bajo coste.

\begin{table}[H]
    \centering
    \caption{Requisitos de rendimiento del sistema}
    \label{tab:req_rendimiento}
    \begin{tabularx}{\linewidth}{l X l}
        \toprule
        \textbf{ID} & \textbf{Descripción} & \textbf{Fuente} \\
        \midrule
        RP01 & Latencia extremo a extremo (sensor a PFD) inferior a \SI{200}{\milli\second}. & \cite{sae_as8034c} \\
        RP02 & Tasa de refresco del PFD igual o superior a \SI{30}{FPS}. & \cite{ac25_11b}, \cite{sae_as8034c} \\
        RP03 & Tasa de actualización de la estimación de actitud igual o superior a \SI{20}{\hertz}. & \cite{ac25_11b} \\
        RP04 & Precisión de la estimación de actitud (cabeceo y alabeo) dentro de $\pm 3^{\circ}$ en condiciones estáticas. & \cite{mahony2008}, \cite{cocoli2023} \\
        RP05 & Precisión de posicionamiento GNSS igual o mejor que \SI{2.5}{\meter} CEP. & \cite{ds_neo6} \\
        RP06 & Alcance del enlace de telemetría suficiente para uso en laboratorio, del orden de decenas de metros en interior. & Decisión de proyecto \\
        \bottomrule
    \end{tabularx}
\end{table}


\section{Restricciones del proyecto}
\label{sec:req_restricciones}

Las restricciones acotan el espacio de diseño y reflejan el contexto educativo y
de bajo coste del proyecto.

\begin{table}[H]
    \centering
    \caption{Restricciones del proyecto}
    \label{tab:req_restricciones}
    \begin{tabularx}{\linewidth}{l X l}
        \toprule
        \textbf{ID} & \textbf{Descripción} & \textbf{Fuente} \\
        \midrule
        RC01 & Coste total de materiales inferior a \SI{100}{\eur}. & Charter \\
        RC02 & Todo el software empleado debe ser de código abierto. & Charter \\
        RC03 & La plataforma debe ser compacta y transportable para uso en laboratorio. & Charter \\
        RC04 & Alimentación mediante conector USB-C. & Charter \\
        RC05 & La PCB debe ser fabricable mediante servicios estándar de prototipado. & Decisión de proyecto \\
        \bottomrule
    \end{tabularx}
\end{table}


\section{Requisitos normativos}
\label{sec:req_normativos}

Los requisitos normativos identifican las normas aeronáuticas que se siguen como
referencia para el diseño. Como se ha indicado al inicio del capítulo, el
sistema no aspira a certificación; las normas se aplican con carácter académico
para que la plataforma sea representativa de un instrumento real.

\begin{table}[H]
    \centering
    \caption{Requisitos normativos del sistema}
    \label{tab:req_normativos}
    \begin{tabularx}{\linewidth}{l X l}
        \toprule
        \textbf{ID} & \textbf{Descripción} & \textbf{Fuente} \\
        \midrule
        RN01 & Disposición de los instrumentos del PFD conforme al patrón Basic-T. & \cite{ac25_11b} \\
        RN02 & Color y simbología del PFD conformes al estándar de displays electrónicos multipropósito. & \cite{sae_as8034c}, \cite{hfstd001b} \\
        RN03 & Jerarquía de color para alertas conforme a la normativa FAA. & \cite{cfr25_1321}, \cite{cfr25_1322} \\
        RN04 & Codificación lógica de la telemetría conforme al protocolo ARINC 429: formato de palabra, etiquetas, codificación BNR, paridad impar y SSM. & \cite{arinc429} \\
        RN05 & El sistema no aspira a certificación aeronáutica ni a uso en vuelo real, por lo que no incorpora redundancia, protección ambiental ni sistemas de seguridad. & Decisión de proyecto, \cref{sec:exclusiones} \\
        \bottomrule
    \end{tabularx}
\end{table}


